\ifdefined\gkver\else\def\gkver{student}\fi
\documentclass[\gkver]{gaokaozhenti}
\begin{document}

\chapter{2023年江苏}

\begin{enumerate}
\item
%% number: 1
%% paperName: 2023年江苏高考第 1 题 4 分
%% typeId: 1
%% type: 单选题
%% chapter: 直线运动
%% point: 运动图象
%% method: 无
%% score: 4
%% degree: 950
%% duplicateId: 0
%% body:
电梯上升过程中，某同学用智能手机记录了电梯速度随时间变化的关系，如图所示。电梯加速上升的时段是\xzanswer{A}

\onepicture[w=7.00cm]{figs/01.png}

\fourchoices[answer=A]
{从 $20.0\Us$ 到 $30.0\Us$}
{从 $30.0\Us$ 到 $40.0\Us$}
{从 $40.0\Us$ 到 $50.0\Us$}
{从 $50.0\Us$ 到 $60.0\Us$}

%% answer: A
%% memo:
\memoanswer{
因电梯上升，由速度图像可知，电梯加速上升的时间段为 $20.0\Us$ 到 $30.0\Us$。

故选 A。
}

\item
%% number: 2
%% paperName: 2023年江苏高考第 2 题 4 分
%% typeId: 1
%% type: 单选题
%% chapter: 磁场
%% point: 安培力
%% method: 无
%% score: 4
%% degree: 750
%% duplicateId: 0
%% body:
如图所示，匀强磁场的磁感应强度为 $B$。L 形导线通以恒定电流 $I$，放置在磁场中。已知 ab 边长为 $2l$，与磁场方向垂直，bc 边长为 $l$，与磁场方向平行。该导线受到的安培力为\xzanswer{C}

\onepicture[w=4.60cm]{figs/02.png}

\fourchoices[answer=C]
{0}
{$BIl$}
{$2BIl$}
{$\sqrt{5}BIl$}

%% answer: C
%% memo:
\memoanswer{
因 bc 段与磁场方向平行，则不受安培力；ab 段与磁场方向垂直，则受安培力为

$F_{ab}=BI\cdot 2l=2BIl$

则该导线受到的安培力为 $2BIl$。

故选 C。
}

\item
%% number: 3
%% paperName: 2023年江苏高考第 3 题 4 分
%% typeId: 1
%% type: 单选题
%% chapter: 气体性质
%% point: 气体图象
%% method: 无
%% score: 4
%% degree: 650
%% duplicateId: 0
%% body:
如图所示，密闭容器内一定质量的理想气体由状态 A 变化到状态 B。该过程中\xzanswer{B}

\onepicture[w=5.00cm]{figs/03.png}

\fourchoices[answer=B]
{气体分子的数密度增大}
{气体分子的平均动能增大}
{单位时间内气体分子对单位面积器壁的作用力减小}
{单位时间内与单位面积器壁碰撞的气体分子数减小}

%% answer: B
%% memo:
\memoanswer{
A．在 $p-T$ 图像中过原点的图线为等容线，即从 A 到 B 气体体积不变，则气体分子的数密度不变，选项 A 错误；

B．从 A 到 B 气体的温度升高，则气体分子的平均动能变大，则选项 B 正确；

C．从 A 到 B 气体的压强变大，气体分子的平均速率变大，则单位时间内气体分子对单位面积的器壁的碰撞力变大，选项 C 错误；

D．气体的分子密度不变，从 A 到 B 气体分子的平均速率变大，则单位时间内与单位面积器壁碰撞的气体分子数变大，选项 D 错误。

故选 B。
}

\item
%% number: 4
%% paperName: 2023年江苏高考第 4 题 4 分
%% typeId: 1
%% type: 单选题
%% chapter: 曲线运动
%% point: 人造地球卫星
%% method: 无
%% score: 4
%% degree: 750
%% duplicateId: 0
%% body:
设想将来发射一颗人造卫星，能在月球绕地球运动的轨道上稳定运行，该轨道可视为圆轨道。该卫星与月球相比，一定相等的是\xzanswer{C}

\fourchoices[answer=C]
{质量}
{向心力大小}
{向心加速度大小}
{受到地球的万有引力大小}

%% answer: C
%% memo:
\memoanswer{
根据 $G\frac{Mm}{r^{2}}=ma$ 可得 $a=\frac{GM}{r^{2}}$。因该卫星与月球的轨道半径相同，可知向心加速度相同；因该卫星的质量与月球质量不同，则向心力大小以及受地球的万有引力大小均不相同。

故选 C。
}

\item
%% number: 5
%% paperName: 2023年江苏高考第 5 题 4 分
%% typeId: 1
%% type: 单选题
%% chapter: 光学
%% point: 光的折射
%% method: 无
%% score: 4
%% degree: 650
%% duplicateId: 0
%% body:
地球表面附近空气的折射率随高度降低而增大，太阳光斜射向地面的过程中会发生弯曲。下列光路图中能描述该现象的是\xzanswer{A}

\fourchoices[answer=A, ispicture=true, h=3.00cm]
{figs/05a.png}
{figs/05b.png}
{figs/05c.png}
{figs/05d.png}

%% answer: A
%% memo:
\memoanswer{
根据折射定律

$n_{\text{上}}\sin\theta_{\text{上}}=n_{\text{下}}\sin\theta_{\text{下}}$

由于地球表面附近空气的折射率随高度降低而增大，则 $n_{\text{下}}>n_{\text{上}}$，则 $\theta_{\text{下}}$ 逐渐减小，画出光路图如下

\onepicture[w=4.60cm, align=right]{figs/05_an.png}

则从高到低 $\theta_{\text{下}}$ 逐渐减小，则光线应逐渐趋于竖直方向。

故选 A。
}

\item
%% number: 6
%% paperName: 2023年江苏高考第 6 题 4 分
%% typeId: 1
%% type: 单选题
%% chapter: 光学
%% point: 光的干涉
%% method: 无
%% score: 4
%% degree: 750
%% duplicateId: 0
%% body:
用某种单色光进行双缝干涉实验，在屏上观察到的干涉条纹如图甲所示，改变双缝间的距离后，干涉条纹如图乙所示，图中虚线是亮纹中心的位置。则双缝间的距离变为原来的\xzanswer{B}

\onepicture[w=7.00cm]{figs/06.png}

\fourchoices[answer=B]
{$\frac{1}{3}$ 倍}
{$\frac{1}{2}$ 倍}
{2 倍}
{3 倍}

%% answer: B
%% memo:
\memoanswer{
根据双缝干涉的条纹间距与波长关系有 $\Delta x=\frac{L}{d}\lambda$。由题图知 $\Delta x_{\text{乙}}=2\Delta x_{\text{甲}}$，则 $d_{\text{乙}}=\frac{1}{2}d_{\text{甲}}$。

故选 B。
}

\item
%% number: 7
%% paperName: 2023年江苏高考第 7 题 4 分
%% typeId: 1
%% type: 单选题
%% chapter: 力物体的平衡
%% point: 多力平衡
%% method: 无
%% score: 4
%% degree: 750
%% duplicateId: 0
%% body:
如图所示，“嫦娥五号”探测器静止在月球平坦表面处。已知探测器质量为 $m$，四条腿与竖直方向的夹角均为 $\theta$，月球表面的重力加速度为地球表面重力加速度 $g$ 的 $\frac{1}{6}$。每条腿对月球表面压力的大小为\xzanswer{D}

\onepicture[w=3.60cm]{figs/07.jpg}

\fourchoices[answer=D]
{$\frac{mg}{4}$}
{$\frac{mg}{4\cos\theta}$}
{$\frac{mg}{6\cos\theta}$}
{$\frac{mg}{24}$}

%% answer: D
%% memo:
\memoanswer{
对“嫦娥五号”探测器受力分析有 $F_{N}=mg_{\text{月}}$，则对一条腿有 $F_{N1}=\frac{1}{4}mg_{\text{月}}=\frac{mg}{24}$。根据牛顿第三定律可知每条腿对月球表面的压力为 $\frac{mg}{24}$。

故选 D。
}

\item
%% number: 8
%% paperName: 2023年江苏高考第 8 题 4 分
%% typeId: 1
%% type: 单选题
%% chapter: 电磁感应
%% point: 右手定则
%% method: 无
%% score: 4
%% degree: 650
%% duplicateId: 0
%% body:
如图所示，圆形区域内有垂直纸面向里的匀强磁场，OC 导体棒的 O 端位于圆心，棒的中点 A 位于磁场区域的边缘。现使导体棒绕 O 点在纸面内逆时针转动。O、A、C 点电势分别为 $\varphi_{O}$、$\varphi_{A}$、$\varphi_{C}$，则\xzanswer{A}

\onepicture[w=5.80cm]{figs/08.png}

\fourchoices[answer=A]
{$\varphi_{O}>\varphi_{C}$}
{$\varphi_{C}>\varphi_{A}$}
{$\varphi_{O}=\varphi_{A}$}
{$\varphi_{O}-\varphi_{A}=\varphi_{A}-\varphi_{C}$}

%% answer: A
%% memo:
\memoanswer{
ABC．由题图可看出 OA 导体棒转动切割磁感线，则根据右手定则可知 $\varphi_{O}>\varphi_{A}$。其中导体棒 AC 段不在磁场中，不切割磁感线，电流为 0，则 $\varphi_{C}=\varphi_{A}$，A 正确、BC 错误；

D．根据以上分析可知 $\varphi_{O}-\varphi_{A}>0$，$\varphi_{A}-\varphi_{C}=0$，则 $\varphi_{O}-\varphi_{A}>\varphi_{A}-\varphi_{C}$。D 错误。

故选 A。
}

\item
%% number: 9
%% paperName: 2023年江苏高考第 9 题 4 分
%% typeId: 1
%% type: 单选题
%% chapter: 气体性质
%% point: 验证玻意耳定律
%% method: 无
%% score: 4
%% degree: 850
%% duplicateId: 0
%% body:
在“探究气体等温变化的规律”的实验中，实验装置如图所示。利用注射器选取一段空气柱为研究对象。下列改变空气柱体积的操作正确的是\xzanswer{B}

\onepicture[w=3.30cm]{figs/09.png}

\fourchoices[answer=B]
{把柱塞快速地向下压}
{把柱塞缓慢地向上拉}
{在橡胶套处接另一注射器，快速推动该注射器柱塞}
{在橡胶套处接另一注射器，缓慢推动该注射器柱塞}

%% answer: B
%% memo:
\memoanswer{
因为该实验是要探究气体等温变化的规律；实验中要缓慢推动或拉动活塞，目的是尽可能保证封闭气体在状态变化过程中的温度不变；为了方便读取封闭气体的体积不需要在橡胶套处接另一注射器。

故选 B。
}

\item
%% number: 10
%% paperName: 2023年江苏高考第 10 题 4 分
%% typeId: 1
%% type: 单选题
%% chapter: 直线运动
%% point: 运动的合成
%% method: 无
%% score: 4
%% degree: 650
%% duplicateId: 0
%% body:
达$\cdot$芬奇的手稿中描述了这样一个实验：一个罐子在空中沿水平直线向右做匀加速运动，沿途连续漏出沙子。若不计空气阻力，则下列图中能反映空中沙子排列的几何图形是\xzanswer{D}

\fourchoices[answer=D, ispicture=true, h=3.00cm]
{figs/10a.png}
{figs/10b.png}
{figs/10c.png}
{figs/10d.png}

%% answer: D
%% memo:
\memoanswer{
罐子在空中沿水平直线向右做匀加速运动，在时间 $\Delta t$ 内水平方向增加量 $a\Delta t^{2}$，竖直方向做自由落体运动，在时间 $\Delta t$ 增加 $g\Delta t^{2}$；说明水平方向位移增加量与竖直方向位移增加量比值一定，则连线的倾角就是一定的。

故选 D。
}

\item
%% number: 11
%% paperName: 2023年江苏高考第 11 题 4 分
%% typeId: 1
%% type: 单选题
%% chapter: 运动和力
%% point: 牛顿定律的应用
%% method: 无
%% score: 4
%% degree: 550
%% duplicateId: 0
%% body:
滑块以一定的初速度沿粗糙斜面从底端上滑，到达最高点 B 后返回到底端。利用频闪仪分别对上滑和下滑过程进行拍摄，频闪照片示意图如图所示。与图乙中相比，图甲中滑块\xzanswer{C}

\onepicture[w=8.00cm]{figs/11.png}

\fourchoices[answer=C]
{受到的合力较小}
{经过 A 点的动能较小}
{在 A、B 之间的运动时间较短}
{在 A、B 之间克服摩擦力做的功较小}

%% answer: C
%% memo:
\memoanswer{
A．因为频闪照片时间间隔相同，对比图甲和乙可知图甲中滑块加速度大，是上滑阶段；根据牛顿第二定律可知图甲中滑块受到的合力较大；故 A 错误；

B．从图甲中的 A 点到图乙中的 A 点，先上升后下降，重力做功为 0，摩擦力做负功；根据动能定理可知图甲经过 A 点的动能较大，故 B 错误；

C．由于图甲中滑块加速度大，根据 $v^{2}-v_{0}^{2}=2ax$ 可知图甲在 A、B 之间的运动时间较短，故 C 正确；

D．由于无论上滑或下滑均受到滑动摩擦力大小相等，故图甲和图乙在 A、B 之间克服摩擦力做的功相等，故 D 错误。

故选 C。
}

\item
%% number: 12
%% paperName: 2023年江苏高考第 12 题 15 分
%% typeId: 4
%% type: 实验题
%% chapter: 电路
%% point: 电路实验
%% method: 无
%% score: 15
%% degree: 450
%% duplicateId: 0
%% body:
小明通过实验探究电压表内阻对测量结果的影响。所用器材有：干电池（电动势约 $1.5\UV$，内阻不计）2 节；两量程电压表（量程 $0\sim 3\UV$，内阻约 $3\UkO$；量程 $0\sim 15\UV$，内阻约 $15\UkO$）1 个；滑动变阻器（最大阻值 $50\UO$）1 个；定值电阻（阻值 $50\UO$）21 个；开关 1 个及导线若干。实验电路如图~\ref{2023江苏12a}所示。
\begin{enumerate}
	\item
	电压表量程应选用\tkanswer{$3\UV$}（选填“3 V”或“15 V”）。
	\item
	该实验的实物电路如图~\ref{2023江苏12b}所示（右侧未拍全）。先将滑动变阻器的滑片置于如图所示的位置，然后用导线将电池盒上接线柱 A 与滑动变阻器的接线柱\tkanswer{D}（选填“B”“C”“D”）连接，再闭合开关，开始实验。
	\item
	将滑动变阻器滑片移动到合适位置后保持不变，依次测量电路中 O 与 1，2，$\ldots$，21 之间的电压。某次测量时，电压表指针位置如图~\ref{2023江苏12c}所示，其示数为\tkanswer{$1.50\UV$}。根据测量数据作出电压 $U$ 与被测电阻值 $R$ 的关系图线，如图~\ref{2023江苏12d}中实线所示。
	\item
	在图~\ref{2023江苏12a}所示的电路中，若电源电动势为 $E$，电压表视为理想电压表，滑动变阻器接入的阻值为 $R_{1}$，定值电阻的总阻值为 $R_{2}$，当被测电阻为 $R$ 时，其两端的电压 $U=$\tkanswer{$\frac{ER}{R_{1}+R_{2}}$}（用 $E$、$R_{1}$、$R_{2}$、$R$ 表示），据此作出 $U-R$ 理论图线如图~\ref{2023江苏12d}中虚线所示。小明发现被测电阻较小或较大时，电压的实测值与理论值相差较小。
	\item
	分析可知，当 $R$ 较小时，$U$ 的实测值与理论值相差较小，是因为电压表的分流小，电压表内阻对测量结果影响较小。小明认为，当 $R$ 较大时，$U$ 的实测值与理论值相差较小，也是因为相同的原因。你是否同意他的观点？请简要说明理由\tkanswer{不同意。当 $R$ 较大时，电压表内阻不能忽略，但因 $R_{V}\gg R_{1}$，电压表分流对测量结果的影响仍较小，故实测值接近理论值}。
\end{enumerate}

\fourpicture[labelA=2023江苏12a, labelB=2023江苏12b, labelC=2023江苏12c, labelD=2023江苏12d, numstyle=arabic, showlabel=false, align=right, wA=7.00cm, wB=7.00cm, wC=4.00cm, wD=5.50cm]
{figs/12a.png}
{figs/12b.png}
{figs/12c.png}
{figs/12d.png}
%% answer:
\jdanswer{
\begin{enumerate}
	\item
	$3\UV$

	\item
	D

	\item
	$1.50\UV$

	\item
	$\frac{ER}{R_{1}+R_{2}}$

	\item
	不同意。当 $R$ 较大时，电压表内阻不能忽略，但因 $R_{V}\gg R_{1}$，电压表分流对测量结果的影响仍较小，故实测值接近理论值
\end{enumerate}
}
%% memo:
\memoanswer{
（1）所用电源为两节干电池，电动势为 $3\UV$，则所用电表量程为 $3\UV$；

（2）闭合电键之前，滑动变阻器阻值应该调到最大，则由图可知，电池盒上的接线柱 A 应该与滑动变阻器的接线柱 D 连接；

（3）电压表最小刻度为 $0.1\UV$，则读数为 $1.50\UV$；

（4）由闭合电路欧姆定律可得 $I=\frac{E}{R_{1}+R_{2}}$

当被测电阻阻值为 $R$ 时电压表读数 $U=IR=\frac{ER}{R_{1}+R_{2}}$

（5）不同意；

当 $R$ 较大时，则电压表内阻不能忽略，则电路中的电流 $I=\frac{E}{R_{1}+(R_{2}-R)+\frac{RR_{V}}{R+R_{V}}}$，则电压表读数为 $U=I\cdot\frac{RR_{V}}{R+R_{V}}=\frac{E}{\frac{(R_{1}+R_{2}-R)(R+R_{V})}{RR_{V}}+1}$，当 $R$ 较大时，$R=R_{2}$ 时 $R$ 最大，此时 $U=\frac{E}{\frac{R_{1}(R_{2}+R_{V})}{R_{2}R_{V}}+1}=\frac{E}{\frac{R_{1}}{R_{V}}+\frac{R_{1}}{R_{2}}+1}$。

因 $R_{V}\gg R_{1}$，则电压表读数接近于 $U=\frac{E}{\frac{R_{1}}{R_{2}}+1}=\frac{ER}{R_{1}+R_{2}}$。
}

\item
%% number: 13
%% paperName: 2023年江苏高考第 13 题 6 分
%% typeId: 6
%% type: 计算题
%% chapter: 曲线运动
%% point: 圆周运动的应用
%% method: 无
%% score: 6
%% degree: 850
%% duplicateId: 0
%% body:
“转碟”是传统的杂技项目，如图所示，质量为 $m$ 的发光物体放在半径为 $r$ 的碟子边缘，杂技演员用杆顶住碟子中心，使发光物体随碟子一起在水平面内绕 A 点做匀速圆周运动。当角速度为 $\omega_{0}$ 时，碟子边缘看似一个光环。求此时发光物体的速度大小 $v_{0}$ 和受到的静摩擦力大小 $f$。

\onepicture[w=3.50cm, align=right]{figs/13.png}
%% answer:
\jdanswer{
$v_{0}=\omega_{0}r$；$f=m\omega_{0}^{2}r$
}
%% memo:
\memoanswer{
发光体的速度 $v_{0}=\omega_{0}r$

发光体做匀速圆周运动，则静摩擦力充当做圆周运动的向心力，则静摩擦力大小为 $f=m\omega_{0}^{2}r$
}

\item
%% number: 14
%% paperName: 2023年江苏高考第 14 题 8 分
%% typeId: 6
%% type: 计算题
%% chapter: 光学
%% point: 光子说
%% method: 建立模型
%% score: 8
%% degree: 650
%% duplicateId: 0
%% body:
“夸父一号”太阳探测卫星可以观测太阳辐射的硬 X 射线。硬 X 射线是波长很短的光子，设波长为 $\lambda$。若太阳均匀地向各个方向辐射硬 X 射线，卫星探测仪镜头正对着太阳，每秒接收到 $N$ 个该种光子。已知探测仪镜头面积为 $S$，卫星离太阳中心的距离为 $R$，普朗克常量为 $h$，光速为 $c$，求：
\begin{enumerate}
	\item
	每个光子的动量 $p$ 和能量 $E$；
	\item
	太阳辐射硬 X 射线的总功率 $P$。
\end{enumerate}
%% answer:
\jdanswer{
\begin{enumerate}
	\item
	$p=\frac{h}{\lambda}$，$E=\frac{hc}{\lambda}$

	\item
	$\frac{4\pi R^{2}Nhc}{S\lambda}$
\end{enumerate}
}
%% memo:
\memoanswer{
（1）由题意可知每个光子的动量为 $p=\frac{h}{\lambda}$

每个光子的能量为 $E=h\nu=\frac{hc}{\lambda}$

（2）太阳均匀地向各个方向辐射硬 X 射线，根据题意设 $t$ 秒发射总光子数为 $n$，则

$\frac{n}{tN}=\frac{4\pi R^{2}}{S}$

可得 $n=\frac{4\pi R^{2}Nt}{S}$

所以 $t$ 秒辐射光子的总能量

$W=E'=nh\frac{c}{\lambda}=\frac{4\pi R^{2}Nthc}{S\lambda}$

太阳辐射硬 X 射线的总功率 $P=\frac{W}{t}=\frac{4\pi R^{2}Nhc}{S\lambda}$
}

\item
%% number: 15
%% paperName: 2023年江苏高考第 15 题 12 分
%% typeId: 6
%% type: 计算题
%% chapter: 曲线运动
%% point: 斜抛运动
%% method: 无
%% score: 12
%% degree: 450
%% duplicateId: 0
%% body:
如图所示，滑雪道 AB 由坡道和水平道组成，且平滑连接，坡道倾角均为 $45^{\circ}$。平台 BC 与缓冲坡 CD 相连。若滑雪者从 P 点由静止开始下滑，恰好到达 B 点。滑雪者现从 A 点由静止开始下滑，从 B 点飞出。已知 A、P 间的距离为 $d$，滑雪者与滑道间的动摩擦因数均为 $\mu$，重力加速度为 $g$，不计空气阻力。
\begin{enumerate}
	\item
	求滑雪者运动到 P 点的时间 $t$；
	\item
	求滑雪者从 B 点飞出的速度大小 $v$；
	\item
	若滑雪者能着陆在缓冲坡 CD 上，求平台 BC 的最大长度 $L$。
\end{enumerate}

\onepicture[w=9.00cm, align=right]{figs/15.png}
%% answer:
\jdanswer{
\begin{enumerate}
	\item
	$t=\sqrt{\frac{2\sqrt{2}d}{g(1-\mu)}}$

	\item
	$\sqrt{\sqrt{2}gd(1-\mu)}$

	\item
	$\sqrt{2}d(1-\mu)$
\end{enumerate}
}
%% memo:
\memoanswer{
（1）滑雪者从 A 到 P 根据动能定理有

$mgd\sin 45^{\circ}-\mu mg\cos 45^{\circ}d=\frac{1}{2}mv_{P}^{2}-0$

根据动量定理有

$\left(mg\sin 45^{\circ}-\mu mg\cos 45^{\circ}\right)t=mv_{P}-0$

联立解得 $t=\sqrt{\frac{2\sqrt{2}d}{g(1-\mu)}}$，$v_{P}=\sqrt{2gd(1-\mu)}$

（2）由于滑雪者从 P 点由静止开始下滑，恰好到达 B 点，故从 P 点到 B 点合力做功为 0，所以当从 A 点下滑时，到达 B 点有

$v_{B}=v_{P}=\sqrt{2gd(1-\mu)}$

（3）当滑雪者刚好落在 C 点时，平台 BC 的长度最大；滑雪者从 B 点飞出做斜抛运动，

竖直方向上有 $v_{P}\cos 45^{\circ}=g\times\frac{t}{2}$

水平方向上有 $L=v_{P}\sin 45^{\circ}\cdot t$

联立可得 $L=\sqrt{2}d(1-\mu)$
}

\item
%% number: 16
%% paperName: 2023年江苏高考第 16 题 15 分
%% typeId: 6
%% type: 计算题
%% chapter: 磁场
%% point: 带电粒子在复合场中的运动
%% method: 无
%% score: 15
%% degree: 450
%% duplicateId: 0
%% body:
霍尔推进器某局部区域可抽象成如图所示的模型。$Oxy$ 平面内存在竖直向下的匀强电场和垂直坐标平面向里的匀强磁场，磁感应强度为 $B$。质量为 $m$、电荷量为 $e$ 的电子从 $O$ 点沿 $x$ 轴正方向水平入射。入射速度为 $v_{0}$ 时，电子沿 $x$ 轴做直线运动；入射速度小于 $v_{0}$ 时，电子的运动轨迹如图中的虚线所示，且在最高点与在最低点所受的合力大小相等。不计重力及电子间相互作用。
\begin{enumerate}
	\item
	求电场强度的大小 $E$；
	\item
	若电子入射速度为 $\frac{v_{0}}{4}$，求运动到速度为 $\frac{v_{0}}{2}$ 时位置的纵坐标 $y_{1}$；
	\item
	若电子入射速度在 $0<v<v_{0}$ 范围内均匀分布，求能到达纵坐标 $y_{2}=\frac{mv_{0}}{5eB}$ 位置的电子数 $N$ 占总电子数 $N_{0}$ 的百分比。
\end{enumerate}

\onepicture[w=8.00cm, align=right]{figs/16.png}
%% answer:
\jdanswer{
\begin{enumerate}
	\item
	$v_{0}B$

	\item
	$\frac{3mv_{0}}{32eB}$

	\item
	90 \%
\end{enumerate}
}
%% memo:
\memoanswer{
（1）由题知，入射速度为 $v_{0}$ 时，电子沿 $x$ 轴做直线运动则有

$Ee=ev_{0}B$

解得 $E=v_{0}B$

（2）电子在竖直向下的匀强电场和垂直坐标平面向里的匀强磁场的复合场中，由于洛伦兹力不做功，且由于电子入射速度为 $\frac{v_{0}}{4}$，则电子受到的电场力大于洛伦兹力，则电子向上偏转，根据动能定理有

$eEy_{1}=\frac{1}{2}m\left(\frac{1}{2}v_{0}\right)^{2}-\frac{1}{2}m\left(\frac{1}{4}v_{0}\right)^{2}$

解得 $y_{1}=\frac{3mv_{0}}{32eB}$

（3）若电子以 $v$ 入射时，设电子能达到的最高点位置的纵坐标为 $y$，则根据动能定理有 $eEy=\frac{1}{2}mv_{m}^{2}-\frac{1}{2}mv^{2}$

由于电子在最高点与在最低点所受的合力大小相等，则在最高点有 $F_{\text{合}}=ev_{m}B-eE$

在最低点有 $F_{\text{合}}=eE-evB$

联立有 $v_{m}=\frac{2E}{B}-v$，$y=\frac{2m(v_{0}-v)}{eB}$

要让电子达纵坐标 $y_{2}=\frac{mv_{0}}{5eB}$ 位置，即 $y\geqslant y_{2}$

解得 $v\leqslant\frac{9}{10}v_{0}$

则若电子入射速度在 $0<v<v_{0}$ 范围内均匀分布，能到达纵坐标 $y_{2}=\frac{mv_{0}}{5eB}$ 位置的电子数 $N$ 占总电子数 $N_{0}$ 的 90\%。
}

\end{enumerate}

\end{document}
